\documentclass[11pt]{article}

\newcommand{\cnum}{Math 245B}
\newcommand{\ced}{Winter 2026}
\newcommand{\ctitle}[4]{\title{\vspace{-0.5in}\cnum, \ced\\Assignment #1: #2}\author{\vspace{-0.35in}\\#3, #4}}
\usepackage{enumitem}
\usepackage[usenames,dvipsnames,svgnames,table,hyperref]{xcolor}
\usepackage{amsmath, amsfonts}
\usepackage{graphicx} % Required for inserting images
\usepackage{float}
\usepackage{listings}
\usepackage{xcolor}
\usepackage{hyperref}
\usepackage{comment}

\renewcommand*{\theenumi}{\alph{enumi}}
\renewcommand*\labelenumi{(\theenumi)}
\renewcommand*{\theenumii}{\roman{enumii}}
\renewcommand*\labelenumii{\theenumii.}

\definecolor{codegreen}{rgb}{0,0.6,0}
\definecolor{codegray}{rgb}{0.5,0.5,0.5}
\definecolor{codepurple}{rgb}{0.58,0,0.82}
\definecolor{backcolour}{rgb}{0.95,0.95,0.92}
\definecolor{header}{HTML}{205459}
\definecolor{solution}{HTML}{204d52}

\newcommand{\solution}[1]{{{\textcolor{header}{Solution:} \textcolor{solution}{#1}}}}
\setlist[enumerate]{label=(\alph*)}

\lstdefinestyle{mystyle}{
    backgroundcolor=\color{backcolour},
    commentstyle=\color{codegreen},
    keywordstyle=\color{magenta},
    numberstyle=\tiny\color{codegray},
    stringstyle=\color{codepurple},
    basicstyle=\ttfamily\footnotesize,
    breakatwhitespace=false,
    breaklines=true,
    captionpos=b,
    keepspaces=true,
    numbers=left,
    numbersep=5pt,
    showspaces=false,
    showstringspaces=false,
    showtabs=false,
    tabsize=2
}


\begin{document}
\ctitle{\#1}{}{Asher Christian}{006-150-286}
\date{}
\maketitle
\section{Wednesday}
Discrete state space, finite or countably infinite continuous time markov property, for any $t_1<t_2<\dots<t_n< t$ and $i,i_1,\dots,i_n,$
$s > 0$
\[
    \mathbb{P}(X_{t+s} = j | X_t = i) = \mathbb{P}(X_{t+s} = j | X_t = i, X_{t_n} = i_n, X_{t_{n-1}} = i_{n-1}, \dots, X_{t_1} = i_1) = p_s(i,j)
\] 
Chapman kolmogorov equations
\[
    P_{t+s}(i,j) = \sum_{k}^{}p_t(i,k)p_s(k,j)
\] 
define transition rate
\[
q(i,j) = \lim_{t\to 0} \frac{p_t(i,j)}{t}
\] 
if it exists. Given rates $(q(i,j), j \ne i)$ you assign independent exponential clock to each transition, each has rate  $q(i,j)$
so you have rate  $q(i,1)$ to jump to state 1 and so on, whichever hits first you go to that state.
The total rate for leaving  $i$ is
\[
\text{exp}(\sum_{j : j \ne i}^{}q(i,j))
\] 
and the probability that the $i-j$ clock ticks is proportional to its rate.
The transition probability from  $i \rightarrow j$ is
\[
\frac{q(i,j)}{\lambda(i)}
\] 
for all $j \ne i$ think of  $X_t$ as embedding in  $ Y_n$
 \[
Y_n : DTMC
\] 
with transition probability $u(i,j) = \frac{q(i,j)}{\lambda_i}$, $j \ne i$. Then  $X_0 = Y_0$ $X_t = Y_0$ for $t < \text{exp}(\lambda(Y_0))$ then
$X_t = Y_1$ for $\text{exp}(\lambda(Y_0) \le t < \text{exp}(\lambda(Y_0)) + \text{exp}(\lambda(Y_1))$ and so on.
Assume $\lambda(i)$ is finite for all $i$. If  $\Lambda = \sup_i \lambda(i) < \infty$ then the rate
for leaving is $\Lambda$. We can construct another embedded discrete time markov chain
with transition probability  $v(i,j) = \frac{q(i.j)}{\Lambda}$ and $v(i,i) = 1 = \frac{\lambda(i)}{\Lambda}$ \\
We use the kolmogorov's backward/forward equations. in the finite state space case
\[
p_t' = Qp_t \qquad p_t' = p_tQ \implies p_t = e^{tQ} = e^{tU D U^{-1}} = U e^{tD} U^{-1}
\] 
if we want $\pi = \pi p_t \implies \pi^{C} Q = 0$ and detailed balance condition:
\[
\pi(i)q(i,j) = \pi(j)q(j,i)
\] 
\[
\pi^{D}(p-I) = 0
\] 
where $\pi^{C}$ is the continuous stationary distribution and $\pi^{D}$ is the stationary distribution
\[
\pi^{C}(i) = \pi^{D}(i) \frac{1}{\lambda(i)}
\] 
limiting behavior: Irreducibility: any $i,j$ there is a sequence of intermediate states  $i=x_0,x_1,\dots,x_n=j$ such that
\[
    q(x_k,x_{k-1}) > 0  \qquad 1 \le k \le n
\] 
thereom if $X_t$ is irreducible and has a stationary distribution  $\pi$ then
\[
\lim_{t\to \infty}p_t(i,j) = \pi(j) \qquad \forall i
\] 
\section{Theorem: Ergodic theorem}
Under the same conditions
\[
\lim_{t\to \infty} \frac{1}{t} \int_{0}^{t}f(X_s)ds = \sum_{i}^{}\pi(i) f(i)
\] 
Exit distribution and Exit time we use the connection between $X_t$ and  $Y_n$ then the exist distribution is the same as  $Y_n$
 \[
h(i) = \mathbb{P}_i(\sigma_A < \sigma_B)
\] 
then
\[
\begin{cases}
    h(i) = 1 & i \in A\\
    h(i) = 0 & i \in B\\
    h(i) = \sum_{j\ne i}^{}u(i,j)h(j) & i \not\in A \cup B
\end{cases}
\] 
Let
\[
Q = 
\begin{pmatrix} 
    A & (A \cup B)^{c} & B\\
      &  & & A\\
    \text{rows that sum to $v$ } & R & & (A \cup B)^{c} \\
      &  & & B
\end{pmatrix} 
\] 
if 
 \[
 R \vec h = -
 \begin{pmatrix} \vdots \\ \sum_{j \in A}^{}Q(i,j) \\ \vdots  \end{pmatrix}  = - \vec v
\] 
for $i \in (A \cup B)^{c}$.\\
For exist time
\[
    g(i) + \mathbb{E}_i[\sigma_A] 
\] 
\[
g(i) = \frac{1}{\lambda_i} + \sum_{j \ne i}^{}u(i,j) g(j) \qquad g(i) = 0, i \in A
\] 
this is equivalent to
\[
R \vec g = -\vec 1 \implies \vec g = -R^{-1}\vec 1
\] 
here
\[
    \begin{pmatrix} & A & A^{c} \\ A & & \\ A^{c} & & R \end{pmatrix} 
\] 
since
\[
    g(i) = \sum_{j \in A^{c}}^{} \mathbb{E}_i[L_j(\sigma_A)]  = -\sum_{j \in A^{c}}^{} R^{-1}(i,j) = -( R^{-1} \vec 1)(i)
\] 
where
\[
L_j(\sigma_A) = \int_{0}^{\sigma_A} 1(X_s = j) ds
\] 



\end{document}
